% ============================================================
%  APPENDIX A: FIGURES
% ============================================================

\section{Figures}
\label{sec:figures}

% ---- Standardized note components ----
% Sample: Tennessee public colleges vs. public colleges in 26
%   non-merit-program states, birth cohorts 1980--1991.
% Estimator: varies per figure.
% SEs: Clustered at the state level (27 clusters).
% CI: 95% confidence intervals.
% Reference: k = -1 omitted.
% Source: Chetty et al. (2020).

% Note: Individual C&S event studies for par_q1 and par_rank have been
% removed. These are subsumed by the four-estimator overlay figure
% (Figure~\ref{fig:four_estimators}) in Section~\ref{sec:robustness_estimators}.

% ============================================================
%  A-0. COHORT 1980 DIAGNOSTIC
% ============================================================

\begin{figure}[H]
  \centering
  \includegraphics[width=0.95\textwidth]{fig_cohort1980_diagnostic.pdf}
  \caption{1980-Cohort Diagnostic: Event Studies With and Without the 1980 Birth Cohort}
  \label{fig:cohort1980_diagnostic}
  \vspace{1mm}
  \begin{minipage}{0.9\textwidth}
    \footnotesize
    \textit{Notes:} TWFE event-study estimates under the preferred
    specification ($t^* = 1987$, national control, BEA baseline county
    income control). Left panels include the 1980 birth cohort ($k = -7$);
    right panels exclude it. Top row: bottom-quintile share ($Q_1$).
    Bottom row: mean parental income rank (\textit{par\_rank}).
    Pre-treatment coefficients at $k = -7$ through $k = -2$ trend away
    from zero when the 1980 cohort is included, indicating a pre-trend
    that violates parallel trends at the longest horizons. Dropping the
    1980 cohort does not materially change the remaining coefficients
    ($k = -6$ through $k = +4$), confirming that the exclusion removes
    an anomalous observation without leveraging the post-treatment estimates.
    $k = -1$ (cohort 1986) is the omitted reference period.
    College and cohort fixed effects. 95\% confidence intervals.
    Standard errors clustered at the state level.
    Data: \textcite{chetty2020mobility}.
  \end{minipage}
\end{figure}

% ============================================================
%  B. QUINTILE DECOMPOSITION
% ============================================================

% Note: Main quintile gradient figure (fig_tn_quintile_gradient) is now
% inline in Section 6.2. Only the event-study version remains here.

\begin{figure}[H]
  \centering
  \includegraphics[width=0.9\textwidth]{fig_quintile_event_studies.pdf}
  \caption{Event-Study Estimates by Income Quintile}
  \label{fig:quintile_event_studies}
  \vspace{1mm}
  \begin{minipage}{0.9\textwidth}
    \footnotesize
    \textit{Notes:} TWFE event-study estimates for each income quintile
    share under the preferred specification ($t^* = 1987$, drop 1980
    cohort, BEA county income control). Bottom quintiles (Q1--Q2) trend
    upward post-treatment (gaining share); top quintiles (Q4--Q5) trend
    downward (losing share); Q3 is near zero. The monotonic gradient
    confirms that HOPE reshuffled composition across the full income
    distribution. National control group (26 non-merit states). $k = -1$
    omitted. Standard errors clustered by state.
    Data: \textcite{chetty2020mobility}.
  \end{minipage}
\end{figure}

% ============================================================
%  C. LEVEL DECOMPOSITION
% ============================================================

\begin{figure}[H]
  \centering
  \includegraphics[width=0.9\textwidth]{fig_share_vs_level.pdf}
  \caption{Share Effects vs.\ Count Effects: Bottom-Quintile Students}
  \label{fig:share_vs_level}
  \vspace{1mm}
  \begin{minipage}{0.9\textwidth}
    \footnotesize
    \textit{Notes:} Comparison of share-based (left) and count-based
    (right) event-study estimates for bottom-quintile students. The share
    increase is accompanied by a count increase, confirming that the
    composition shift reflects genuine sorting rather than a denominator
    artifact. Count outcomes are inherently noisier than share outcomes
    because they inherit enrollment variation across colleges (coefficient
    of variation $\approx 1.0$). The wider confidence intervals on count
    estimates reflect this additional noise, not weaker identification.
    TWFE estimator. Data: \textcite{chetty2020mobility}.
  \end{minipage}
\end{figure}

\begin{figure}[H]
  \centering
  \includegraphics[width=0.9\textwidth]{fig_level_decomposition.pdf}
  \caption{Level Decomposition: Bottom-Quintile Count, Top-Quintile Count, and Total Enrollment}
  \label{fig:level_decomp}
  \vspace{1mm}
  \begin{minipage}{0.9\textwidth}
    \footnotesize
    \textit{Notes:} TWFE event-study estimates for count outcomes (quintile share $\times$ cohort
    enrollment). The critical finding is that total enrollment (Panel~B)
    is unchanged post-treatment, ruling out the denominator-artifact
    explanation for share-based results. Q1 count increases and Q5 count
    decreases, consistent with real compositional sorting. Count outcomes
    exhibit wider confidence intervals than share outcomes because they
    compound share-level variation with enrollment variation (enrollment
    coefficient of variation $\approx 1.0$ across colleges). The
    pre-treatment noise in count event studies reflects this inherent
    volatility rather than a failure of parallel trends for the
    share-based specifications. Data: \textcite{chetty2020mobility}.
  \end{minipage}
\end{figure}

% ============================================================
%  D. MECHANISM EVIDENCE
% ============================================================

\begin{figure}[H]
  \centering
  \includegraphics[width=0.9\textwidth]{fig_tn_privates_spillover.pdf}
  \caption{Spillover Test: Tennessee Private Colleges}
  \label{fig:spillover}
  \vspace{1mm}
  \begin{minipage}{0.9\textwidth}
    \footnotesize
    \textit{Notes:} TWFE event-study estimates comparing Tennessee
    private colleges (treated) to private
    colleges in control states. No significant composition change is
    detected, indicating that top-quintile students who exit Tennessee
    public colleges do not re-sort to in-state private institutions.
    $N = 5{,}336$. Data: \textcite{chetty2020mobility}.
  \end{minipage}
\end{figure}

\begin{figure}[H]
  \centering
  \includegraphics[width=0.9\textwidth]{fig_within_tn_comparison.pdf}
  \caption{Within-Tennessee Comparison: Public vs.\ Private Colleges}
  \label{fig:within_tn}
  \vspace{1mm}
  \begin{minipage}{0.9\textwidth}
    \footnotesize
    \textit{Notes:} Event-study estimates restricting the sample to
    Tennessee institutions only. Public colleges (treated) are compared
    to Tennessee private colleges (control). The composition shift is
    concentrated in the public sector. College and cohort fixed effects
    included. $N = 571$. Data: \textcite{chetty2020mobility}.
  \end{minipage}
\end{figure}

% Note: Estimator comparisons (LP-DiD, TWFE, Sun-Abraham, dCDH) are now
% consolidated in the four-estimator overlay figure (Figure~\ref{fig:four_estimators})
% in Section~\ref{sec:robustness_estimators}. Individual TWFE event studies
% removed to avoid redundancy.

% ============================================================
%  G. SENSITIVITY ANALYSIS
% ============================================================

% Note: R&R sensitivity is now presented as a 2×2 panel figure
% (fig_rr_combined) inline in Section 7 (Robustness). The old
% breakdown-by-event-time figure has been removed; the combined
% figure provides a cleaner summary.

% ============================================================
%  H. SPATIAL HETEROGENEITY
% ============================================================

\begin{figure}[H]
  \centering
  \includegraphics[width=0.9\textwidth]{fig_spatial_scatter.pdf}
  \caption{College-Level DiD Effects vs.\ Baseline County Income}
  \label{fig:spatial_scatter}
  \vspace{1mm}
  \begin{minipage}{0.9\textwidth}
    \footnotesize
    \textit{Notes:} Each point represents a Tennessee public college.
    The horizontal axis is baseline log per capita income (BEA county
    average, 1998--2003); the vertical axis is the college-specific
    TWFE DiD estimate for bottom-quintile share. Separate OLS fitted
    lines and 95\% confidence bands are shown for metro (navy) and
    nonmetro (cranberry) colleges. The overall correlation is
    $r = 0.36$; richer counties host colleges with larger progressive
    effects. Data: \textcite{chetty2020mobility}; BEA CAINC1.
  \end{minipage}
\end{figure}

\begin{figure}[H]
  \centering
  \includegraphics[width=0.9\textwidth]{fig_subgroup_forest.pdf}
  \caption{Subgroup ATT Estimates: Metro/Nonmetro $\times$ Sector and Tier}
  \label{fig:subgroup_forest}
  \vspace{1mm}
  \begin{minipage}{0.9\textwidth}
    \footnotesize
    \textit{Notes:} TWFE static DiD estimates for bottom-quintile share,
    estimated separately by metro status, sector, and tier. Metro public
    colleges---especially two-year institutions---drive the aggregate
    effect. Nonmetro public colleges show small, insignificant effects.
    Private colleges show an attenuated metro pattern consistent with
    HOPE's smaller private-institution subsidy. Horizontal bars are 95\%
    confidence intervals. Standard errors clustered at the state level.
    Data: \textcite{chetty2020mobility}.
  \end{minipage}
\end{figure}

\begin{figure}[H]
  \centering
  \includegraphics[width=0.9\textwidth]{fig_multistate_metro.pdf}
  \caption{Metro Interaction Coefficients Across States}
  \label{fig:multistate_metro}
  \vspace{1mm}
  \begin{minipage}{0.9\textwidth}
    \footnotesize
    \textit{Notes:} Metro $\times$ treatment interaction coefficients
    from TWFE static DiD, estimated separately for each state. Open
    circles show overall ATT for context. Tennessee, South Dakota, and
    Georgia all show positive metro interactions, indicating that the
    metro amplification is not specific to Tennessee. Massachusetts
    shows no metro variation because virtually all public colleges are
    in metro areas. Horizontal bars are 95\% confidence intervals.
    Standard errors clustered at the state level.
    Data: \textcite{chetty2020mobility}.
  \end{minipage}
\end{figure}

% ============================================================
%  I. AWARD STRUCTURE AND STATE RANKINGS
% ============================================================

\begin{figure}[H]
  \centering
  \includegraphics[width=0.85\textwidth]{fig_state_buffering.pdf}
  \caption{Economic Resilience by Program Design}
  \label{fig:state_buffering}
  \vspace{1mm}
  \begin{minipage}{0.85\textwidth}
    \footnotesize
    \textit{Notes:} Shock $\times$ merit-state interaction coefficient on
    bottom-quintile share for each strong merit program (Sjoquist \&
    Winters, 2015, plus Massachusetts). Each state estimated separately
    against non-merit baseline. Public and private nonprofit colleges.
    BEA county income shock at decision year (cohort~$+$~17). College
    and cohort FE, college-clustered SEs. 95\% CIs.
  \end{minipage}
\end{figure}
